\section{Modules of differentials}
\label{section-differentials}

\noindent
In this section we briefly explain how to define the module of relative
differentials for a morphism of ringed spaces. We suggest the reader take
a look at the corresponding section in the chapter on commutative algebra
(Algebra, Section \ref{algebra-section-differentials}).

\begin{definition}
\label{definition-derivation}
Let $X$ be a topological space. Let $\varphi : \mathcal{O}_1 \to \mathcal{O}_2$
be a homomorphism of sheaves of rings. Let $\mathcal{F}$ be an
$\mathcal{O}_2$-module. An {\it $\mathcal{O}_1$-derivation} or more precisely
a {\it $\varphi$-derivation} into $\mathcal{F}$ is a map
$D : \mathcal{O}_2 \to \mathcal{F}$ which is additive, annihilates the image
of $\mathcal{O}_1 \to \mathcal{O}_2$, and satisfies the
{\it Leibniz rule}
$$
D(ab) = aD(b) + D(a)b
$$
for all $a, b$ local sections of $\mathcal{O}_2$ (wherever they are both
defined). We denote $\text{Der}_{\mathcal{O}_1}(\mathcal{O}_2, \mathcal{F})$
the set of $\varphi$-derivations into $\mathcal{F}$.
\end{definition}

\noindent
This is the sheaf theoretic analogue of
Algebra, Definition \ref{algebra-definition-derivation}.
Given a derivation $D : \mathcal{O}_2 \to \mathcal{F}$
as in the definition the map on global sections
$$
D : \Gamma(X, \mathcal{O}_2) \longrightarrow \Gamma(X, \mathcal{F})
$$
is a $\Gamma(X, \mathcal{O}_1)$-derivation as in the algebra definition.
Note that if $\alpha : \mathcal{F} \to \mathcal{G}$ is a map of
$\mathcal{O}_2$-modules, then there is an induced map
$$
\text{Der}_{\mathcal{O}_1}(\mathcal{O}_2, \mathcal{F})
\longrightarrow
\text{Der}_{\mathcal{O}_1}(\mathcal{O}_2, \mathcal{G})
$$
given by the rule $D \mapsto \alpha \circ D$. In other words
we obtain a functor.

\begin{lemma}
\label{lemma-universal-module}
Let $X$ be a topological space. Let $\varphi : \mathcal{O}_1 \to \mathcal{O}_2$
be a homomorphism of sheaves of rings. The functor
$$
\textit{Mod}(\mathcal{O}_2) \longrightarrow \textit{Ab}, \quad
\mathcal{F} \longmapsto \text{Der}_{\mathcal{O}_1}(\mathcal{O}_2, \mathcal{F})
$$
is representable.
\end{lemma}

\begin{proof}
This is proved in exactly the same way as the analogous statement in algebra.
During this proof, for any sheaf of sets $\mathcal{F}$ on $X$,
let us denote $\mathcal{O}_2[\mathcal{F}]$ the sheafification of the
presheaf $U \mapsto \mathcal{O}_2(U)[\mathcal{F}(U)]$ where this denotes
the free $\mathcal{O}_2(U)$-module on the set $\mathcal{F}(U)$.
For $s \in \mathcal{F}(U)$ we denote $[s]$ the corresponding section
of $\mathcal{O}_2[\mathcal{F}]$ over $U$. If $\mathcal{F}$ is a sheaf of
$\mathcal{O}_2$-modules, then there is a canonical map
$$
c : \mathcal{O}_2[\mathcal{F}] \longrightarrow \mathcal{F}
$$
which on the presheaf level is given by the rule
$\sum f_s[s] \mapsto \sum f_s s$. We will employ the short hand
$[s] \mapsto s$ to
describe this map and similarly for other maps below. Consider
the map of $\mathcal{O}_2$-modules
\begin{equation}
\label{equation-define-module-differentials}
\begin{matrix}
\mathcal{O}_2[\mathcal{O}_2 \times \mathcal{O}_2] \oplus
\mathcal{O}_2[\mathcal{O}_2 \times \mathcal{O}_2] \oplus
\mathcal{O}_2[\mathcal{O}_1] &
\longrightarrow &
\mathcal{O}_2[\mathcal{O}_2] \\
[(a, b)] \oplus [(f, g)] \oplus [h] & \longmapsto & [a + b] - [a] - [b] + \\
& & [fg] - g[f] - f[g] + \\
& & [\varphi(h)]
\end{matrix}
\end{equation}
with short hand notation as above. Set $\Omega_{\mathcal{O}_2/\mathcal{O}_1}$
equal to the cokernel of this map. Then it is clear that there exists
a map of sheaves of sets
$$
\text{d} : \mathcal{O}_2 \longrightarrow \Omega_{\mathcal{O}_2/\mathcal{O}_1}
$$
mapping a local section $f$ to the image of $[f]$ in
$\Omega_{\mathcal{O}_2/\mathcal{O}_1}$. By construction $\text{d}$
is a $\mathcal{O}_1$-derivation. Next, let $\mathcal{F}$
be a sheaf of $\mathcal{O}_2$-modules and let
$D : \mathcal{O}_2 \to \mathcal{F}$ be a $\mathcal{O}_1$-derivation.
Then we can consider the $\mathcal{O}_2$-linear map
$\mathcal{O}_2[\mathcal{O}_2] \to \mathcal{F}$ which sends $[g]$ to $D(g)$.
It follows from the definition of a derivation that this map annihilates
sections in the image of the map (\ref{equation-define-module-differentials})
and hence defines a map
$$
\alpha_D : \Omega_{\mathcal{O}_2/\mathcal{O}_1} \longrightarrow \mathcal{F}
$$
Since it is clear that $D = \alpha_D \circ \text{d}$ the lemma is proved.
\end{proof}

\begin{definition}
\label{definition-module-differentials}
Let $X$ be a topological space. Let $\varphi : \mathcal{O}_1 \to \mathcal{O}_2$
be a homomorphism of sheaves of rings on $X$. The {\it module of differentials}
of $\varphi$ is the object representing the functor
$\mathcal{F} \mapsto \text{Der}_{\mathcal{O}_1}(\mathcal{O}_2, \mathcal{F})$
which exists by Lemma \ref{lemma-universal-module}.
It is denoted $\Omega_{\mathcal{O}_2/\mathcal{O}_1}$, and the {\it universal
$\varphi$-derivation} is denoted
$\text{d} : \mathcal{O}_2 \to \Omega_{\mathcal{O}_2/\mathcal{O}_1}$.
\end{definition}

\noindent
Note that $\Omega_{\mathcal{O}_2/\mathcal{O}_1}$ is the cokernel of
the map (\ref{equation-define-module-differentials}) of
$\mathcal{O}_2$-modules. Moreover the map $\text{d}$ is described
by the rule that $\text{d}f$ is the image of the local section $[f]$.

\begin{lemma}
\label{lemma-differentials-sheafify}
Let $X$ be a topological space. Let $\varphi : \mathcal{O}_1 \to \mathcal{O}_2$
be a homomorphism of sheaves of rings on $X$. Then
$\Omega_{\mathcal{O}_2/\mathcal{O}_1}$ is the sheaf associated to the
presheaf $U \mapsto \Omega_{\mathcal{O}_2(U)/\mathcal{O}_1(U)}$.
\end{lemma}

\begin{proof}
Consider the map (\ref{equation-define-module-differentials}). There is
a similar map of presheaves whose value on the open $U$ is
$$
\mathcal{O}_2(U)[\mathcal{O}_2(U) \times \mathcal{O}_2(U)] \oplus
\mathcal{O}_2(U)[\mathcal{O}_2(U) \times \mathcal{O}_2(U)] \oplus
\mathcal{O}_2(U)[\mathcal{O}_1(U)]
\longrightarrow
\mathcal{O}_2(U)[\mathcal{O}_2(U)]
$$
The cokernel of this map has value $\Omega_{\mathcal{O}_2(U)/\mathcal{O}_1(U)}$
over $U$ by the construction of the module of differentials in 
Algebra, Definition \ref{algebra-definition-differentials}.
On the other hand, the sheaves in (\ref{equation-define-module-differentials})
are the sheafifications of the presheaves above. Thus the result follows
as sheafification is exact.
\end{proof}

\begin{lemma}
\label{lemma-localize-differentials}
Let $X$ be a topological space. Let $\varphi : \mathcal{O}_1 \to \mathcal{O}_2$
be a homomorphism of sheaves of rings. For $U \subset X$ open
there is a canonical isomorphism
$$
\Omega_{\mathcal{O}_2/\mathcal{O}_1}|_U =
\Omega_{(\mathcal{O}_2|_U)/(\mathcal{O}_1|_U)}
$$
compatible with universal derivations.
\end{lemma}

\begin{proof}
Holds because $\Omega_{\mathcal{O}_2/\mathcal{O}_1}$
is the cokernel of the map (\ref{equation-define-module-differentials}).
\end{proof}

\begin{lemma}
\label{lemma-pullback-differentials}
Let $f : Y \to X$ be a continuous map of topological spaces.
Let $\varphi : \mathcal{O}_1 \to \mathcal{O}_2$
be a homomorphism of sheaves of rings on $X$.
Then there is a canonical identification
$f^{-1}\Omega_{\mathcal{O}_2/\mathcal{O}_1} =
\Omega_{f^{-1}\mathcal{O}_2/f^{-1}\mathcal{O}_1}$
compatible with universal derivations.
\end{lemma}

\begin{proof}
This holds because the sheaf $\Omega_{\mathcal{O}_2/\mathcal{O}_1}$
is the cokernel of the map (\ref{equation-define-module-differentials})
and a similar statement holds for
$\Omega_{f^{-1}\mathcal{O}_2/f^{-1}\mathcal{O}_1}$,
because the functor $f^{-1}$ is exact, and because
$f^{-1}(\mathcal{O}_2[\mathcal{O}_2]) =
f^{-1}\mathcal{O}_2[f^{-1}\mathcal{O}_2]$,
$f^{-1}(\mathcal{O}_2[\mathcal{O}_2 \times \mathcal{O}_2]) =
f^{-1}\mathcal{O}_2[f^{-1}\mathcal{O}_2 \times f^{-1}\mathcal{O}_2]$, and
$f^{-1}(\mathcal{O}_2[\mathcal{O}_1]) =
f^{-1}\mathcal{O}_2[f^{-1}\mathcal{O}_1]$.
\end{proof}

\begin{lemma}
\label{lemma-stalk-module-differentials}
Let $X$ be a topological space. Let $\mathcal{O}_1 \to \mathcal{O}_2$
be a homomorphism of sheaves of rings on $X$. Let $x \in X$. Then we have
$\Omega_{\mathcal{O}_2/\mathcal{O}_1, x} =
\Omega_{\mathcal{O}_{2, x}/\mathcal{O}_{1, x}}$.
\end{lemma}

\begin{proof}
This is a special case of Lemma \ref{lemma-pullback-differentials}
for the inclusion map $\{x\} \to X$. An alternative proof is to use
Lemma \ref{lemma-differentials-sheafify},
Sheaves, Lemma \ref{sheaves-lemma-stalk-sheafification}, and
Algebra, Lemma \ref{algebra-lemma-colimit-differentials}.
\end{proof}

\begin{lemma}
\label{lemma-functoriality-differentials}
Let $X$ be a topological space. Let
$$
\xymatrix{
\mathcal{O}_2 \ar[r]_\varphi & \mathcal{O}_2' \\
\mathcal{O}_1 \ar[r] \ar[u] & \mathcal{O}'_1 \ar[u]
}
$$
be a commutative diagram of sheaves of rings on $X$. The map
$\mathcal{O}_2 \to \mathcal{O}'_2$ composed with the map
$\text{d} : \mathcal{O}'_2 \to \Omega_{\mathcal{O}'_2/\mathcal{O}'_1}$
is a $\mathcal{O}_1$-derivation. Hence we obtain a canonical map of
$\mathcal{O}_2$-modules
$\Omega_{\mathcal{O}_2/\mathcal{O}_1} \to
\Omega_{\mathcal{O}'_2/\mathcal{O}'_1}$.
It is uniquely characterized by the property that
$\text{d}(f) \mapsto \text{d}(\varphi(f))$
for any local section $f$ of $\mathcal{O}_2$.
In this way $\Omega_{-/-}$ becomes a functor on the category
of arrows of sheaves of rings.
\end{lemma}

\begin{proof}
This lemma proves itself.
\end{proof}

\begin{lemma}
\label{lemma-differential-seq}
In Lemma \ref{lemma-functoriality-differentials} suppose that
$\mathcal{O}_2 \to \mathcal{O}'_2$ is surjective with kernel
$\mathcal{I} \subset \mathcal{O}_2$ and assume that
$\mathcal{O}_1 = \mathcal{O}'_1$. Then there is a canonical exact
sequence of $\mathcal{O}'_2$-modules
$$
\mathcal{I}/\mathcal{I}^2
\longrightarrow
\Omega_{\mathcal{O}_2/\mathcal{O}_1} \otimes_{\mathcal{O}_2} \mathcal{O}'_2
\longrightarrow
\Omega_{\mathcal{O}'_2/\mathcal{O}_1}
\longrightarrow
0
$$
The leftmost map is characterized by the rule that a local section
$f$ of $\mathcal{I}$ maps to $\text{d}f \otimes 1$.
\end{lemma}

\begin{proof}
For a local section $f$ of $\mathcal{I}$ denote $\overline{f}$ the image of
$f$ in $\mathcal{I}/\mathcal{I}^2$. To show that the map
$\overline{f} \mapsto \text{d}f \otimes 1$ is well defined we just have to
check that $\text{d} f_1f_2 \otimes 1 = 0$ if $f_1, f_2$ are local sections
of $\mathcal{I}$. And this is clear from the Leibniz rule
$\text{d} f_1f_2 \otimes 1 =
(f_1 \text{d}f_2 + f_2 \text{d} f_1 )\otimes 1 =
\text{d}f_2 \otimes f_1 + \text{d}f_1 \otimes f_2 = 0$.
A similar computation show this map is
$\mathcal{O}'_2 = \mathcal{O}_2/\mathcal{I}$-linear. The map on the right
is the one from Lemma \ref{lemma-functoriality-differentials}. To see
that the sequence is exact, we can check on stalks
(Lemma \ref{lemma-abelian}). By
Lemma \ref{lemma-stalk-module-differentials}
this follows from
Algebra, Lemma \ref{algebra-lemma-differential-seq}.
\end{proof}

\begin{definition}
\label{definition-differentials}
Let $(f, f^\sharp) : (X, \mathcal{O}_X) \to (S, \mathcal{O}_S)$
be a morphism of ringed spaces.
\begin{enumerate}
\item Let $\mathcal{F}$ be an $\mathcal{O}_X$-module. An {\it $S$-derivation}
into $\mathcal{F}$ is a $f^{-1}\mathcal{O}_S$-derivation, or more
precisely a $f^\sharp$-derivation in the sense of
Definition \ref{definition-derivation}.
We denote $\text{Der}_S(\mathcal{O}_X, \mathcal{F})$
the set of $S$-derivations into $\mathcal{F}$.
\item The {\it sheaf of differentials $\Omega_{X/S}$ of $X$ over $S$}
is the module of differentials $\Omega_{\mathcal{O}_X/f^{-1}\mathcal{O}_S}$
endowed with its universal
$S$-derivation $\text{d}_{X/S} : \mathcal{O}_X \to \Omega_{X/S}$.
\end{enumerate}
\end{definition}

\noindent
Here is a particular situation where derivations come up
naturally.

\begin{lemma}
\label{lemma-double-structure-gives-derivation}
Let $(f, f^\sharp) : (X, \mathcal{O}_X) \to (S, \mathcal{O}_S)$
be a morphism of ringed spaces. Consider a short exact sequence
$$
0 \to \mathcal{I} \to \mathcal{A} \to \mathcal{O}_X \to 0
$$
Here $\mathcal{A}$ is a sheaf of $f^{-1}\mathcal{O}_S$-algebras,
$\pi : \mathcal{A} \to \mathcal{O}_X$ is a surjection
of sheaves of $f^{-1}\mathcal{O}_S$-algebras, and
$\mathcal{I} = \Ker(\pi)$ is its kernel. Assume $\mathcal{I}$ an ideal
sheaf with square zero in $\mathcal{A}$. So $\mathcal{I}$
has a natural structure of an $\mathcal{O}_X$-module.
A section $s : \mathcal{O}_X \to \mathcal{A}$ of $\pi$
is a $f^{-1}\mathcal{O}_S$-algebra map such that $\pi \circ s = \text{id}$.
Given any section $s : \mathcal{O}_X \to \mathcal{A}$
of $\pi$ and any $S$-derivation $D : \mathcal{O}_X \to \mathcal{I}$
the map
$$
s + D : \mathcal{O}_X \to \mathcal{A}
$$
is a section of $\pi$ and every section $s'$ is of the form $s + D$
for a unique $S$-derivation $D$.
\end{lemma}

\begin{proof}
Recall that the $\mathcal{O}_X$-module structure on $\mathcal{I}$
is given by $h \tau = \tilde h \tau$ (multiplication in $\mathcal{A}$)
where $h$ is a local section of $\mathcal{O}_X$, and
$\tilde h$ is a local lift of $h$ to a local
section of $\mathcal{A}$, and $\tau$ is a local section of $\mathcal{I}$.
In particular, given $s$, we may use $\tilde h = s(h)$.
To verify that $s + D$ is a homomorphism of sheaves of rings we
compute
\begin{eqnarray*}
(s + D)(ab) & = & s(ab) + D(ab) \\
& = & s(a)s(b) + aD(b) + D(a)b \\
& = & s(a) s(b) + s(a)D(b) + D(a)s(b) \\
& = & (s(a) + D(a))(s(b) + D(b))
\end{eqnarray*}
by the Leibniz rule. In the same manner one shows
$s + D$ is a $f^{-1}\mathcal{O}_S$-algebra
map because $D$ is an $S$-derivation. Conversely, given $s'$ we set
$D = s' - s$. Details omitted.
\end{proof}

\begin{lemma}
\label{lemma-functoriality-differentials-ringed-spaces}
Let
$$
\xymatrix{
X' \ar[d]_{h'} \ar[r]_f & X \ar[d]^h \\
S' \ar[r]^g & S
}
$$
be a commutative diagram of ringed spaces.
\begin{enumerate}
\item The canonical map $\mathcal{O}_X \to f_*\mathcal{O}_{X'}$ composed with
$f_*\text{d}_{X'/S'} : f_*\mathcal{O}_{X'} \to f_*\Omega_{X'/S'}$ is a
$S$-derivation and we obtain a canonical map of $\mathcal{O}_X$-modules
$\Omega_{X/S} \to f_*\Omega_{X'/S'}$.
\item The commutative diagram
$$
\xymatrix{
f^{-1}\mathcal{O}_X \ar[r] & \mathcal{O}_{X'} \\
f^{-1}h^{-1}\mathcal{O}_S \ar[u] \ar[r] & (h')^{-1}\mathcal{O}_{S'} \ar[u]
}
$$
induces by Lemmas \ref{lemma-pullback-differentials} and
\ref{lemma-functoriality-differentials}
a canonical map $f^{-1}\Omega_{X/S} \to \Omega_{X'/S'}$.
\end{enumerate}
These two maps correspond (via adjointness of $f_*$ and $f^*$ and
via $f^*\Omega_{X/S} =
f^{-1}\Omega_{X/S} \otimes_{f^{-1}\mathcal{O}_X} \mathcal{O}_{X'}$ and
Sheaves, Lemma \ref{sheaves-lemma-adjointness-tensor-restrict})
to the same $\mathcal{O}_{X'}$-module homomorphism
$$
c_f : f^*\Omega_{X/S} \longrightarrow \Omega_{X'/S'}
$$
which is uniquely characterized by the property that
$f^*\text{d}_{X/S}(a)$ maps to $\text{d}_{X'/S'}(f^*a)$
for any local section $a$ of $\mathcal{O}_X$.
\end{lemma}

\begin{proof}
Omitted.
\end{proof}

\begin{lemma}
\label{lemma-check-functoriality-differentials}
Let
$$
\xymatrix{
X'' \ar[d] \ar[r]_g & X' \ar[d] \ar[r]_f & X \ar[d] \\
S'' \ar[r] & S' \ar[r] & S
}
$$
be a commutative diagram of ringed spaces. With notation as in
Lemma \ref{lemma-functoriality-differentials-ringed-spaces} we have
$$
c_{f \circ g} = c_g \circ g^* c_f
$$
as maps $(f \circ g)^*\Omega_{X/S} \to \Omega_{X''/S''}$.
\end{lemma}

\begin{proof}
Omitted.
\end{proof}

\begin{lemma}
\label{lemma-triangle-differentials}
Let $f : X \to Y$, $g : Y \to S$ be morphisms of ringed spaces
Then there is a canonical exact sequence
$$
f^*\Omega_{Y/S} \to \Omega_{X/S} \to \Omega_{X/Y} \to 0
$$
where the maps come from applications of
Lemma \ref{lemma-functoriality-differentials-ringed-spaces}.
\end{lemma}

\begin{proof}
By taking induced maps in stalks at $x \in X$ and using
Lemma \ref{lemma-stalk-module-differentials} we obtain a sequence
$$
\mathcal{O}_{X, x} \otimes_{\mathcal{O}_{Y, f(x)}}
\Omega_{\mathcal{O}_{Y, f(x)}/\mathcal{O}_{S, g(f(x))}} \to
\Omega_{\mathcal{O}_{X, x}/\mathcal{O}_{S, g(f(x))}} \to
\Omega_{\mathcal{O}_{X, x}/\mathcal{O}_{Y, f(x)}} \to 0
$$
It suffices to see that the maps of the sequence are the same as
the ones in Algebra, Lemma \ref{algebra-lemma-exact-sequence-differentials}.
This is because the characterization of the maps in
Lemma \ref{lemma-functoriality-differentials-ringed-spaces}
and because via the isomorphism of
Sheaves, Lemma \ref{sheaves-lemma-stalk-pullback-modules}
we have $(f^*s)_x = s_x \otimes 1$,
for any local section $s$ of a sheaf of $\mathcal{O}_Y$-modules.
\end{proof}


















